\justifying
\NSFCBodyText

\subsubsection{研究方法}

对货物或运输批次 \(i\)、模态 \(m\)、样本 \(j\) 和事件 \(k\)，记事件时间为 \(\tau\)，采样时间为 \(s_{imj}\)，到达时间为 \(a_{imj}\)，观测与质量量为 \(O_{imj},Q_{imj}\)，链路与资源状态为 \(N_i(t),B_i(t)\)。在同一在线信息边界下，分别定义运输回溯检测风险和仓储前瞻预测风险：
\[
\begin{aligned}
\mathcal F_{i,t}&=\sigma\!\left(\{s_{imj},a_{imj},R_{imj},O_{imj},Q_{imj}:a_{imj}\le t\},N_i([0,t]),B_i([0,t])\right),\\
\theta^{\rm det}_{ik}(t,h_-)&=\Pr\!\left(Y_{ik}(\tau)=1,\tau\in(t-h_-,t]\mid\mathcal F_{i,t}\right),\\
r^{\rm pred}_{ik}(t,h_+)&=\Pr\!\left(Y_{ik}(\tau)=1,\tau\in(t,t+h_+]\mid\mathcal F_{i,t}\right).
\end{aligned}
\tag{1}
\]
式中 \(R_{imj}\) 为采样指示。C1估计已经发生但证据异步到达的运输事件，C2降低这类关键事件在截止期前的漏警风险，C3预测仓储未来风险；三者的在线输出均只能是 \(\mathcal F_{i,t}\) 的函数，截止期后补传数据和事后标签只用于评分。

\textbf{运输事件风险辨识。} 以低速常开记录、设备健康、链路状态、采样指令及队列/重传日志组成可观测锚点 \(W\)，将潜在事件、采样、观测和到达过程分解为
\[
\begin{aligned}
p(X,Y,O,R,A\mid W,Q,N,B)
&=p(X,Y\mid W)\prod_{m,j}p(R_{mj}\mid\mathcal H_s,W,Q,N,B)\\
&\quad\times p(O_{mj}\mid X(s_{mj}),Y,Q)\\
&\quad\times p(A_{mj}\mid R_{mj},O_{mj},\mathcal H_s,W,Q,N,B).
\end{aligned}
\tag{2}
\]
其中 \(X,Y\) 为状态和事件，\(R,A\) 分别表示是否形成样本和何时送达。在锚点覆盖的支持域内，检验给定历史后的序贯可忽略性近似和正值性；满足条件时估计点风险。若事件状态对采样/到达仍有未观测依赖，则用选择优势比敏感性参数 \(\Gamma\ge1\) 约束偏离程度，给出部分辨识区间 \([L^\Gamma_{ik,t},U^\Gamma_{ik,t}]\)；支持不足时拒判。训练目标为
\[
\mathcal L_1=\sum_{i,k,t}w_k(y_{ik,t}-\hat\theta^{\rm det}_{ik,t})^2-\lambda_o\log p_\psi(R,A\mid\mathcal H,W,Q,N,B)
+\lambda_c\Omega_{\rm cov}([L^\Gamma_{ik,t},U^\Gamma_{ik,t}],y_{ik,t}).
\tag{3}
\]
覆盖项用于检验区间而不替代辨识假设。通过去采样时钟、去到达时钟、合并未采/未达及破坏锚点的压力试验，联合评价点风险校准、区间覆盖、区间宽度和拒判率。

\textbf{事件证据价值与安全决策。} 令 \(E_t\) 为已到达证据集合，\(I_t\) 包含缓存摘要、网络、能量和截止期余量，\(a_t\) 为采样/压缩、缓存/传输或边缘计算动作，\(\Delta E_D(a_t)\) 为截止期 \(D\) 前实际新增且可解码的证据。冻结事件判别器 \(\delta\)，以单实体条件期望损失 \(\rho_D(E,I)=\mathbb E[\ell_D(Y,\delta(E))\mid E,I]\) 定义价值，并以置信下界（lower confidence bound, LCB）抑制高不确定动作：
\[
V_t(a_t\mid E_t,I_t)=\rho_D(E_t,I_t)-\mathbb E[\rho_D(E_t\cup\Delta E_D(a_t),I_{t+1})\mid I_t,a_t],\quad
V_t^{\rm LCB}=\widehat V_\phi-\kappa\widehat\sigma_\phi.
\tag{4}
\]
LCB用价值估计减去不确定性惩罚。受控试验先以高可靠参考采集形成带完整时间戳的母流，再以共同随机数重放候选采样、压缩、链路、缓存和续传动作，使未执行动作的潜在证据及资源状态可追溯；真实单路径日志只在行为策略支持重叠域内用于离策略估计，支持不足时不给动作外推。预测器、价值模型和策略按上层实体交叉拟合，最终测试标签封存并只开启一次，不参与在线动作或价值监督。

安全屏障先过滤会破坏端侧最低告警、关键原始证据保留、压缩保真或幂等续传的动作，再求解
\[
\min_\pi\ \mathbb E_\pi[\mathcal R_D(E_D)]+\lambda_B\mathbb E_\pi[C_{\rm byte}]+\lambda_P\mathbb E_\pi[C_{\rm energy}],\quad
\text{s.t. }C_{\rm byte}\le\bar B,\ C_{\rm energy}\le\bar P,\ a_t\in\mathcal A_{\rm safe}.
\tag{5}
\]
式中 \(\pi\) 为策略，\(\bar B,\bar P\) 为通信和能量预算。证据不可静默丢弃、分片校验一致和续传状态幂等属于协议层硬不变量；极端断连且无可行动作时转入端侧最低告警或拒判。识别漏警、误警与超期属于统计风险终点，不承诺在所有链路条件下为零。

\textbf{仓储风险校准与条件性事件记忆。} H3a以温湿度、烟雾、门禁和视频 \(Z_i^w(t)\) 建立独立的仓储本地风险；H3b只在货物级配对和统计功效达到进入条件后检验运输事件记忆：
\[
\begin{aligned}
H_i^w(t)&=F_\eta(H_i^w(t^-),Z_i^w(t),Q_i^w(t)),\quad
r^{\rm pred}_{ik,0}(t,h_+)=g_{k,0}(H_i^w(t),C_i),\\
U_{ik}(t)&=\sum_{e\in\mathcal E_{ik}^{\rm tr}(t_0)}w_{ik,e}\exp[-\lambda_k(t-t_e)],\quad
r^{\rm pred}_{ik}(t,h_+)=g_k(H_i^w(t),U_{ik}(t),C_i),\quad t\ge t_0.
\end{aligned}
\tag{6}
\]
运输事件集合 \(\mathcal E_{ik}^{\rm tr}(t_0)\) 在入库时冻结，记忆在仓储阶段继续衰减；比较无衰减、共享衰减和事件特异衰减。H3b仅保留冲击/倾斜—包装损伤、非法开启—货物安全等合理链条，烟雾、外部侵入和仓储温湿度风险由H3a判断。校准风险与不确定性映射为记录、人工复核和处置建议三级响应，并用漏触发、误触发和触发时延评价；支持域外或关键证据不足时拒判。

\subsubsection{技术路线}

项目采用“一套数据结构、一个边缘原型、三项递进实验”。事件、采样、到达、质量、缓存/重传和标签规范构成统一数据底座，可控异常与参数可追溯网络回放形成母流。三时钟风险及辨识区间用于驱动集合价值、安全采样/压缩、离线缓存、断点续传、可靠上传和边缘计算。仓储本地实时监测、漂移校准和分级响应构成必做闭环，事件记忆只在配对门通过后检验。各内容共享数据规范、基线、评估脚本和原型，每种模态不再另建独立算法。

\subsubsection{实验手段}

数据按货物/运输批次、车辆、路段、仓库和时间分组，同一事件窗口、近重复视频帧、增强及网络扰动副本不得跨集合。训练、校准/阈值选择和最终测试严格分离。C1与具有相同时间戳、质量变量和容量的缺失掩码连续时间模型比较，并做 \(\Gamma\) 敏感性分析；C2用同一母流、随机数、动作、安全屏障和预算，与逐包可加截止期调度比较；C3的H3a比较规则阈值、后处理校准、漂移再校准和同覆盖率选择性预测，H3b再比较真实、分层置换和时间错配运输摘要。统一采用组平衡时间依赖Brier分数（group-balanced Brier score, GBS）和截止期加权漏警风险：
\[
\operatorname{GBS}=\frac1G\sum_g\frac1{|\mathcal T_g|}\sum_{t\in\mathcal T_g}(y_{g,t}-\hat p_{g,t})^2,\quad
\mathcal R_D=\frac1G\sum_g\frac{\sum_{o\in\mathcal O_g}c_{k(o)}\mathbb I(o\text{ 在 }D_{k(o)}\text{ 前漏报})}{\sum_{o\in\mathcal O_g}c_{k(o)}}.
\tag{7}
\]
GBS避免长轨迹支配，\(\mathcal R_D\) 按每个上层实体的异常机会归一化。以最上层独立实体自助重采样报告区间；另报告校准斜率/可靠性图、选择性风险—覆盖曲线、固定召回下误报、证据链交付率、通信字节、估算或实测能耗、缓存溢出和P95/P99时延。效应量、非劣界、样本量、截止期和主次终点层级均由不进入正式测试集的预实验与功效分析冻结。

\subsubsection{关键技术}

关键技术包括：在可观测锚点和敏感性分析下统一三时钟风险辨识、区间与拒判；用可重放母流和支持域诊断估计集合条件价值，并使硬安全屏障先于资源优化；在H3a独立校准基础上，以门控配对、负对照和组外测试检验H3b事件记忆。

\subsubsection{可行性分析}

申请人在时间序列、长时序预测、多源信息融合、多模态视觉理解、状态空间建模、机器学习和软件研发方面具有积累，可支撑风险建模、价值学习和原型实现；新疆大学相关平台与人工智能算力中心可支撑训练、仿真和数据处理。现有材料不能证明现场准入、传感/边缘设备、真实2G/3G轨迹、物理功耗测量或货物级运输—仓储配对已具备，因此先在基础可控场景中建立母流和参数化弱网模拟，再根据设备与授权条件进入代表性板卡和有限场景验证。没有实测时只报告仿真或估算结果；配对门不通过时只完成H3a，不宣称H3b得到验证。
